\chapter{Linearity: changing the language of a problem}\label{ch:linear}
\goal{Work with vectors as ordered lists of numbers, and recognize a linear map as a map that preserves addition and scaling. See why a matrix, a finite table of numbers, can describe such a map completely, and why matrices are multiplied the way they are. Use the kernel to decide when a linear map is injective, and prove the Cauchy--Schwarz and triangle inequalities. Throughout, notice how a change of language, such as a new basis or a matrix in place of a formula, can explain why a fact is true instead of only confirming it by calculation.}

\section{Vectors are ordered measurements}\label{ch10:sec:vectors}
Let $n$ be a positive integer. A \emph{vector} in $\RR^n$ is an ordered list $x=(x_1,x_2,\ldots,x_n)$ of $n$ real numbers, and the numbers $x_1,\ldots,x_n$ are its \emph{coordinates}. So $\RR^n$ is the set of all such lists. For $n=2$ the vectors are ordered pairs $(x_1,x_2)$ of real numbers. We picture $(x_1,x_2)$ as the point of the plane that lies $x_1$ units to the right of a fixed point, the \emph{origin}, and $x_2$ units above it. The coordinates need not describe a position, though. A shop could record $3$ kilograms of flour and $5$ kilograms of sugar as the vector $(3,5)$. Order matters: as with the ordered pairs of Section~\ref{ch03:sec:sets}, $(1,3)$ and $(3,1)$ are different vectors, and for the shop they describe different stock. Two vectors are equal exactly when they agree in every coordinate.

Vectors are added coordinate by coordinate. To multiply a vector by a real number $c$, multiply every coordinate by $c$:
\begin{gather*}
(x_1,\ldots,x_n)+(y_1,\ldots,y_n)=(x_1+y_1,\ldots,x_n+y_n),\\
c(x_1,\ldots,x_n)=(cx_1,\ldots,cx_n).
\end{gather*}
For example, $(1,3)+(2,-1)=(3,2)$ and $2(1,3)=(2,6)$. In this setting a real number such as $c$ is called a \emph{scalar}, because multiplying by it rescales the vector. The \emph{zero vector} has every coordinate equal to zero. We write it simply as $0$ and let the context tell us whether $0$ means a number or a vector. The \emph{negative} of $x$ is $-x=(-1)x=(-x_1,\ldots,-x_n)$, and subtraction means adding the negative: $x-y=x+(-y)=(x_1-y_1,\ldots,x_n-y_n)$. Because every coordinate is handled separately, the familiar rules of arithmetic carry over from numbers to vectors. For instance, $x+y=y+x$, $x+0=x$, $x-x=0$, and $c(x+y)=cx+cy$, because each of these equations holds in every coordinate.

A \emph{linear combination} of vectors $v_1,\ldots,v_k$ is a vector of the form
\[
a_1v_1+a_2v_2+\cdots+a_kv_k,
\]
where $a_1,\ldots,a_k$ are scalars, called the \emph{coefficients} of the combination. Here the subscripts on $v_1,\ldots,v_k$ number the vectors in a list; each $v_j$ is a whole vector, not a coordinate. For example,
\[
2(1,3)-(2,-1)=(2,6)+(-2,1)=(0,7)
\]
is a linear combination of $(1,3)$ and $(2,-1)$ with coefficients $2$ and $-1$.

In the plane, the vectors $e_1=(1,0)$ and $e_2=(0,1)$ serve as building blocks. Every vector is a linear combination of them, because
\[
x_1e_1+x_2e_2=(x_1,0)+(0,x_2)=(x_1,x_2).
\]
Moreover, the coefficients are forced. If $(x_1,x_2)=a_1e_1+a_2e_2$, the same calculation shows that the right side is $(a_1,a_2)$. Two ordered pairs are equal only when their entries agree, so $a_1=x_1$ and $a_2=x_2$.

A list of vectors in $\RR^n$ is called a \emph{basis} of $\RR^n$ when every vector of $\RR^n$ can be written as a linear combination of the list in exactly one way. We have just shown that $e_1,e_2$ is a basis of $\RR^2$. The same argument works in $\RR^n$ for the vectors $e_1,\ldots,e_n$, where $e_j$ has a $1$ in position $j$ and a $0$ in every other position. This list is called the \emph{standard basis} of $\RR^n$. In $\RR^3$, for instance, $(4,5,0)=4e_1+5e_2+0e_3$.

\subsection*{Two requirements: reaching every vector, and reaching it only once}
The definition of a basis makes two separate demands. First, every vector must have a representation, as $(3,-2)=3e_1-2e_2$ does in the standard basis. A list \emph{spans} $\RR^n$ if every vector of $\RR^n$ is a linear combination of the list. Second, the representation must be \emph{unique}: the same vector must never arise from two different lists of coefficients. A basis is a list that does both.

Neither demand implies the other. The single vector $(1,1)$ passes the uniqueness test: a vector that it does make, it makes in only one way, because $a(1,1)=b(1,1)$ forces $a=b$. But it does not span the plane, because every multiple $a(1,1)=(a,a)$ has equal coordinates, so no multiple equals $(1,0)$. In the other direction, the list $(1,0),(0,1),(1,1)$ does span the plane, since already
\[
(x_1,x_2)=x_1(1,0)+x_2(0,1)+0(1,1).
\]
But it does not give unique coefficients. The vector $(1,1)$ can be made as $0(1,0)+0(0,1)+1(1,1)$, and also as $1(1,0)+1(0,1)+0(1,1)$.

Uniqueness has a convenient test. A list $v_1,\ldots,v_k$ is called \emph{linearly independent} if the only linear combination of it that equals the zero vector is the one whose coefficients are all zero:
\[
a_1v_1+\cdots+a_kv_k=0\quad\text{implies}\quad a_1=a_2=\cdots=a_k=0.
\]
A list is linearly independent exactly when no vector can be written as a combination of the list in two different ways. To see why, suppose first that the list is linearly independent and that some vector $x$ has two representations,
\[
x=a_1v_1+\cdots+a_kv_k\qquad\text{and}\qquad x=b_1v_1+\cdots+b_kv_k.
\]
Subtract the second from the first and collect the terms that contain each $v_j$:
\[
(a_1-b_1)v_1+\cdots+(a_k-b_k)v_k=x-x=0.
\]
By independence, every difference $a_j-b_j$ is zero, so the two representations were the same after all. Conversely, if the list is not linearly independent, then some combination with at least one nonzero coefficient equals $0$. The combination with all coefficients zero also equals $0$, so the zero vector has two different representations.

For the list $(1,0),(0,1),(1,1)$, subtracting the second representation of $(1,1)$ found above from the first gives
\[
(-1)(1,0)+(-1)(0,1)+1(1,1)=(0,0),
\]
a combination equal to zero whose coefficients are not all zero. So this list is not linearly independent, which is another way of saying that it does not give unique coefficients. In short, a basis is a list that spans and is linearly independent.

\pauseq{Is the list $(1,2),(2,4)$ a basis of $\RR^2$?}{No, and it fails both requirements. It is not linearly independent, because $2(1,2)+(-1)(2,4)=(0,0)$. It does not span the plane either: every combination $a(1,2)+b(2,4)=(a+2b,\,2a+4b)$ has second coordinate equal to twice its first, so $(1,0)$ is never reached.}

\subsection*{A basis can be chosen to fit the problem}
The standard basis is not the only basis of the plane. Take $w_1=(1,1)$ and $w_2=(1,-1)$. Which coefficients $a$ and $b$ give $aw_1+bw_2=(x_1,x_2)$? Since $aw_1+bw_2=(a+b,\,a-b)$, comparing coordinates shows that we need
\[
a+b=x_1\qquad\text{and}\qquad a-b=x_2.
\]
Adding these two equations gives $2a=x_1+x_2$, and subtracting the second from the first gives $2b=x_1-x_2$. So the only possible coefficients are
\[
a=\frac{x_1+x_2}{2}\qquad\text{and}\qquad b=\frac{x_1-x_2}{2},
\]
and these coefficients do work, since with them $a+b=x_1$ and $a-b=x_2$. Every vector of the plane therefore has exactly one representation, and $w_1,w_2$ is a basis. For instance, $(3,1)=2w_1+1w_2$. In this basis, the coefficients of a vector are half the sum and half the difference of its usual coordinates. The vectors themselves have not changed; only the way we describe them has. Choosing a description that suits the question is a theme of this chapter, and this particular basis will return several times.

\section{A linear map is determined by its building blocks}
A map $T:\RR^n\to\RR^m$ is linear when it preserves addition and scalar multiplication. Equivalently, for all vectors $x,y$ and real scalars $a,b$,
\[
T(ax+by)=aT(x)+bT(y).
\]
For a linear map on the plane,
\[
T(x_1,x_2)=x_1T(e_1)+x_2T(e_2).
\]
Knowing the two outputs $T(e_1),T(e_2)$ therefore tells us every output. This is a mathematical explanation of why a finite array can describe a transformation acting on infinitely many inputs.

Suppose $T(e_1)=(1,1)$ and $T(e_2)=(1,-1)$. Then
\[
T(x_1,x_2)=(x_1+x_2,x_1-x_2).
\]
We record the two output vectors as columns of a matrix:
\[
A=\begin{pmatrix}1&1\\1&-1\end{pmatrix},\qquad
Ax=\begin{pmatrix}x_1+x_2\\x_1-x_2\end{pmatrix}.
\]
An $m$-by-$n$ matrix has $m$ rows and $n$ columns. Its action is defined by
\[
(Ax)_i=\sum_{j=1}^n A_{ij}x_j\qquad(1\le i\le m).
\]
The subscript $ij$ identifies the entry in row $i$, column $j$. The formula says to multiply entries in matching positions of a row and the input list, and then add those products. We will call this operation a dot product below.

For example, if a row is $(2,-1,3)$ and the input is $(4,5,0)$, the corresponding output coordinate is
\[
 2\cdot4+(-1)\cdot5+3\cdot0=3.
\]
The input must have three coordinates because that row has three entries. An $m$-by-$n$ matrix takes $n$-coordinate inputs and produces $m$-coordinate outputs. The two letters indicate possibly different sizes; a matrix need not be square.

We sometimes write a vector horizontally as $(x_1,x_2)$ and sometimes vertically as a column. They name the same ordered data here. In a product $Ax$, we use the column convention so the displayed array agrees with the row-by-input formula.

\pauseq{Is $F(x_1,x_2)=(x_1+1,x_2)$ linear?}{No. Every linear map must send zero to zero: $T(0)=T(0+0)=T(0)+T(0)$ implies $T(0)=0$. Here $F(0,0)=(1,0)$. A translation adds a fixed vector to every input. It preserves differences between inputs but generally moves the zero vector. Such a rule is an example of an affine map; we need only the concrete distinction from a linear map here.}

\section{Kernel: the information lost at zero}
The kernel of a linear map is the set of inputs sent to zero:
\[
\ker T=\{x:T(x)=0\}.
\]
If two inputs have the same output, their difference belongs to the kernel because
\[
T(x-y)=T(x)-T(y)=0.
\]
Thus every collision can be translated into an invisible nonzero input.

\par\noindent\begin{minipage}{\linewidth}
\begin{theorem}\label{thm:kernel}
A linear map is injective if and only if its kernel consists only of the zero vector.
\end{theorem}
\end{minipage}\par

\begin{proof}
If $T$ is injective and $T(x)=0$, then $T(x)=T(0)$. Injectivity gives $x=0$, so the kernel contains no other vector.

Conversely, suppose the only vector mapped to zero is zero. If $T(x)=T(y)$, linearity gives $T(x-y)=0$. The kernel assumption forces $x-y=0$, hence $x=y$. Therefore $T$ is injective.
\end{proof}

The kernel theorem is a two-direction statement about a \emph{linear} map. In its first direction, linearity ensures $T(0)=0$, so an input with output zero can be compared with the zero input. In its second direction, linearity is needed to turn a collision into $T(x-y)=0$. For a nonlinear map, inspecting only inputs sent to zero may miss collisions elsewhere. For instance, $F(x)=x^2$ sends only zero to zero, yet $F(1)=F(-1)$.

This also separates solving a particular equation $T(x)=b$ from proving all targets attainable. A zero kernel says there cannot be two different solutions for any one target. It does not assert that a solution exists for every target. The inclusion $t\mapsto(t,0)$ from $\RR$ to $\RR^2$ is linear and injective, but no input maps to $(0,1)$. To see why counting equations alone cannot establish uniqueness, consider $T(x_1,x_2)=(x_1+x_2,2x_1+2x_2)$. It has two outputs but loses the nonzero vector $(1,-1)$. The second equation contains no new distinction beyond the first.

For the earlier matrix $A$, apply its sum-and-difference rule twice:
\begin{align*}
 A(Ax)&=A(x_1+x_2,x_1-x_2)\\
 &=\big((x_1+x_2)+(x_1-x_2),\ (x_1+x_2)-(x_1-x_2)\big)\\
 &=(2x_1,2x_2)=2x.
\end{align*}
Therefore $Ax=0$ implies $2x=0$, hence $x=0$. The map is injective. Moreover, for any target $b$, taking $x=Ab/2$ gives $Ax=A(Ab)/2=b$. This proves surjectivity and explicitly constructs the inverse. We have proved reversibility without introducing determinants.

\section{The geometry of an inner product}
For vectors $u,v\in\RR^n$, define their dot product by
\[
u\cdot v=\sum_{j=1}^n u_jv_j.
\]
The dot product is also called the standard \emph{inner product} on $\RR^n$; the two names refer to the same operation.

Symmetry follows because each product $u_jv_j$ equals $v_ju_j$. The phrase ``linear in each argument'' means that one argument can be held fixed while addition and scaling in the other distribute through the dot product. For example,
\[
 (a u+b w)\cdot v
 =\sum_j(au_j+bw_j)v_j
 =a\sum_j u_jv_j+b\sum_j w_jv_j
 =a(u\cdot v)+b(w\cdot v).
\]
The same expansion works in the second argument. All sums here contain exactly $n$ terms. The Euclidean length is
\[
\norm{u}_2=\sqrt{u\cdot u}=\sqrt{\sum_{j=1}^n u_j^2}.
\]
Each square is nonnegative, so the square root is defined. Length is zero exactly when every coordinate is zero.

The subscript $2$ on $\norm{u}_2$ identifies this square-based length. It is not the second coordinate of $u$ and does not restrict the vector to two coordinates. Vectors are orthogonal when their dot product is zero. In the plane this agrees with perpendicularity, but the algebraic definition works in any finite dimension.

For $u=(1,1)$ and $v=(1,-1)$, $u\cdot v=1-1=0$. Orthogonality says the positive and negative contributions cancel. This same cancellation will express zero correlation between sign sequences in Chapter~14.

\section{Deriving a bound by removing a component}\label{ch10:sec:cs}
How large can a dot product be, compared with the lengths of the two vectors? A natural guess is
\[
|u\cdot v|\le\norm{u}_2\norm{v}_2.
\]
Let us test it on $u=(3,1)$ and $v=(1,1)$. Here $u\cdot v=4$, $\norm{u}_2=\sqrt{10}$ and $\norm{v}_2=\sqrt2$, so the guess claims that $4\le\sqrt{10}\sqrt2$. To check this without decimals, compare squares. For nonnegative numbers $p$ and $q$, we have $p\le q$ exactly when $p^2\le q^2$. Section~\ref{sec:magnitude} showed that $0\le p\le q$ implies $p^2\le q^2$. Conversely, if $p>q\ge0$, the strict version of the same rule gives $p^2>q^2$, so $p^2\le q^2$ rules out $p>q$. Here the squares are $16$ and $10\cdot2=20$, and $16\le20$, so the guess holds in this example.

How could we find a reason that works for all vectors? Start from a fact that holds for every real number $t$: the vector $u-tv$ has a nonnegative squared length,
\[
(u-tv)\cdot(u-tv)\ge0.
\]
Each choice of $t$ gives a different inequality, and we want an informative one. A good choice removes from $u$ its part along $v$, so that what is left over, $u-tv$, is orthogonal to $v$. So we solve
\[
(u-tv)\cdot v=0.
\]
By linearity in the first argument, this equation says $u\cdot v-t(v\cdot v)=0$. When $v\ne0$, positivity gives $v\cdot v>0$, so we may divide by $v\cdot v$ and obtain
\[
t=\frac{u\cdot v}{v\cdot v}.
\]
This is the value of $t$ used in the proof below. With this $t$, the vector $tv$ is called the \emph{component} of $u$ along $v$, and $u-tv$ is the \emph{residual}: what remains of $u$ after its component along $v$ is removed. In our example $t=4/2=2$, so the component is $2v=(2,2)$ and the residual is $u-2v=(3,1)-(2,2)=(1,-1)$. The residual has dot product $1-1=0$ with $v=(1,1)$, as planned, and its squared length is $2$. The picture below shows this example: the residual meets the dashed line through $v$ at a right angle.

\begin{center}
\begin{tikzpicture}[scale=1.15,>=Stealth]
\draw[->,gray] (-0.3,0)--(3.8,0);
\draw[->,gray] (0,-0.3)--(0,2.8);
\draw[dashed,gray] (-0.3,-0.3)--(2.7,2.7);
\draw[->,thick] (0,0)--(3,1) node[right]{$u=(3,1)$};
\draw[->] (0,0)--(2,2) node[above left]{$2v=(2,2)$};
\draw[->,very thick] (0,0)--(1,1) node[above left]{$v$};
\draw[->,thick] (2,2)--(3,1) node[midway,above right]{$u-2v$};
\draw (1.85,1.85)--(2,1.7)--(2.15,1.85);
\end{tikzpicture}
\end{center}

The theorem below bounds the absolute value $|u\cdot v|$. That is more informative than a bound on $u\cdot v$ alone, because a dot product can be negative. By Section~\ref{sec:magnitude}, the inequality $|u\cdot v|\le\norm{u}_2\norm{v}_2$ says that $u\cdot v$ lies between $-\norm{u}_2\norm{v}_2$ and $\norm{u}_2\norm{v}_2$.

\par\noindent\begin{minipage}{\linewidth}
\begin{theorem}[Cauchy--Schwarz inequality]\label{thm:cs}
For all $u,v\in\RR^n$, $|u\cdot v|\le\norm{u}_2\norm{v}_2$.
\end{theorem}
\end{minipage}\par

\noindent\textit{Plan.} The squared length of $u-tv$ is nonnegative for every real $t$. Expand it, and then substitute the value of $t$ found above.
\begin{proof}
If $v=0$, then $u\cdot v=0$ and $\norm{v}_2=0$, so both sides of the inequality are zero and it holds. Assume from now on that $v\ne0$, so that $v\cdot v>0$. For any real $t$, expand the squared length of $u-tv$, using linearity first in the first argument and then in the second:
\begin{align*}
 (u-tv)\cdot(u-tv)
 &=u\cdot(u-tv)-t\bigl(v\cdot(u-tv)\bigr)\\
 &=u\cdot u-t(u\cdot v)-t(v\cdot u)+t^2(v\cdot v)\\
 &=u\cdot u-2t(u\cdot v)+t^2(v\cdot v).
\end{align*}
The last step uses symmetry, $v\cdot u=u\cdot v$. The left side is a squared length, so by positivity the right side is nonnegative for every real $t$.

Now choose $t=(u\cdot v)/(v\cdot v)$, which is defined because $v\cdot v>0$. Then $t(u\cdot v)=(u\cdot v)^2/(v\cdot v)$, and
\[
t^2(v\cdot v)=\frac{(u\cdot v)^2}{(v\cdot v)^2}\,(v\cdot v)=\frac{(u\cdot v)^2}{v\cdot v}.
\]
Substituting these two expressions gives
\begin{align*}
0&\le u\cdot u-2\,\frac{(u\cdot v)^2}{v\cdot v}+\frac{(u\cdot v)^2}{v\cdot v}\\
&=u\cdot u-\frac{(u\cdot v)^2}{v\cdot v}.
\end{align*}
Add the fraction to both sides and multiply by the positive number $v\cdot v$:
\[
(u\cdot v)^2\le(u\cdot u)(v\cdot v)=\norm{u}_2^2\,\norm{v}_2^2.
\]
Finally, $|u\cdot v|$ and $\norm{u}_2\norm{v}_2$ are nonnegative numbers whose squares are $(u\cdot v)^2$ and $\norm{u}_2^2\norm{v}_2^2$, because $|a|^2=a^2$ for every real number $a$. Nonnegative numbers compare in the same way as their squares, so $|u\cdot v|\le\norm{u}_2\norm{v}_2$.
\end{proof}

\subsection*{When does equality hold?}
When is $|u\cdot v|$ exactly equal to $\norm{u}_2\norm{v}_2$? For $v\ne0$, the answer is: exactly when $u$ is a scalar multiple of $v$.

The proof already contains the reason. With the chosen $t$, it showed that the squared length of the residual is
\[
\norm{u-tv}_2^2=u\cdot u-\frac{(u\cdot v)^2}{v\cdot v}.
\]
Multiplying by $v\cdot v$ gives
\[
 (u\cdot u)(v\cdot v)-(u\cdot v)^2
 =(v\cdot v)\,\norm{u-tv}_2^2.
\]
The left side is the gap between the two sides of the squared inequality. Two nonnegative numbers are equal exactly when their squares are equal, so equality holds in the theorem exactly when this gap is zero. Since $v\cdot v>0$, the gap is zero exactly when $\norm{u-tv}_2=0$, that is, when the residual $u-tv$ is the zero vector, which means $u=tv$. So if equality holds, then $u$ is a multiple of $v$. Conversely, suppose that $u=sv$ for some number $s$. Then
\[
t=\frac{(sv)\cdot v}{v\cdot v}=\frac{s(v\cdot v)}{v\cdot v}=s,
\]
so the residual is $u-sv=0$, the gap is zero, and equality holds.

For example, $u=(2,2)$ and $v=(1,1)$ give $u\cdot v=4$ and $\norm{u}_2\norm{v}_2=\sqrt8\sqrt2=2\sqrt2\cdot\sqrt2=4$. This is equality, as expected, since $u=2v$. For $u=(3,1)$ and $v=(1,1)$ the gap is $20-16=4$, and indeed $(v\cdot v)\norm{u-2v}_2^2=2\cdot2=4$. When $v=0$, both sides of the inequality are zero, so equality holds for every $u$. Altogether, equality holds exactly when $v=0$ or $u$ is a scalar multiple of $v$.

Section~\ref{ch09:sec:taste} recommended asking, of every bound, what happens when it is attained. Here the answer tells us two things. First, the bound is as good as possible: no constant factor smaller than $1$ can be placed in front of its right-hand side. When $u=v\ne0$, both sides equal $\norm{v}_2^2$, so a smaller right-hand side, such as $0.9\,\norm{u}_2\norm{v}_2$, would fail for these vectors. Second, the gap measures how far $u$ is from being a multiple of $v$: it equals $v\cdot v$ times the squared length of the residual.

\subsection*{From the dot-product bound to a distance bound}
In the plane, walking from one point to another by way of a third point is never shorter than walking there directly. Chapter~\ref{ch:limits} needs this fact in $\RR^n$, and Cauchy--Schwarz gives it quickly. For any $u,v\in\RR^n$, expand using linearity in each argument and symmetry:
\begin{align*}
 \norm{u+v}_2^2
 &=(u+v)\cdot(u+v)\\
 &=u\cdot u+u\cdot v+v\cdot u+v\cdot v\\
 &=\norm{u}_2^2+2(u\cdot v)+\norm{v}_2^2\\
 &\le\norm{u}_2^2+2|u\cdot v|+\norm{v}_2^2\\
 &\le\norm{u}_2^2+2\norm{u}_2\norm{v}_2+\norm{v}_2^2\\
 &=(\norm{u}_2+\norm{v}_2)^2.
\end{align*}
The first inequality uses $a\le|a|$ for every real number $a$. The second uses Cauchy--Schwarz. The last line is the expansion $(p+q)^2=p^2+2pq+q^2$. Both $\norm{u+v}_2$ and $\norm{u}_2+\norm{v}_2$ are nonnegative, and nonnegative numbers compare in the same way as their squares. We obtain the \emph{triangle inequality}
\[
 \norm{u+v}_2\le\norm{u}_2+\norm{v}_2.
\]
To see where the name comes from, take three points $x,y,z\in\RR^n$ and put $u=x-y$ and $v=y-z$. Then $u+v=x-z$, and the inequality becomes
\[
\norm{x-z}_2\le\norm{x-y}_2+\norm{y-z}_2.
\]
In the plane, $\norm{x-z}_2$ is the ordinary distance between the points $x$ and $z$, by Pythagoras' theorem. So the inequality says that in the triangle with corners $x$, $y$, $z$, the side from $x$ to $z$ is at most as long as the other two sides together. Exercise~\ref{ex:10.5} asks you to rebuild this argument on your own.

\section{Matrix multiplication records composition}
If $A$ maps $\RR^n$ to $\RR^m$ and $B$ maps $\RR^m$ to $\RR^k$, their product $BA$ records ``$B$ after $A$.'' Expanding one coordinate gives
\[
(B(Ax))_i=\sum_{r=1}^m B_{ir}\sum_{j=1}^n A_{rj}x_j
=\sum_{j=1}^n\left(\sum_{r=1}^m B_{ir}A_{rj}\right)x_j.
\]
The sums are finite, so regrouping is legitimate. This derives $(BA)_{ij}=\sum_rB_{ir}A_{rj}$. Matrix multiplication is not an arbitrary convention; it is forced by composition of the actions.

To inspect the regrouping without nested sums, take both $A$ and $B$ to be two-by-two matrices. The first coordinate after applying $B$ to $Ax$ has the form
\begin{align*}
 B_{11}(A_{11}x_1+A_{12}x_2)
 &+B_{12}(A_{21}x_1+A_{22}x_2)\\
 &=(B_{11}A_{11}+B_{12}A_{21})x_1\\
 &\quad +(B_{11}A_{12}+B_{12}A_{22})x_2.
\end{align*}
The coefficients multiplying $x_1$ and $x_2$ are the first row of the composed matrix. In the general formula, the index $r$ runs through the intermediate coordinates, while $j$ runs through the original inputs. These are different jobs. The summation notation is a compact record of this same distributive calculation, not a new operation.

The transpose $A^{\mathsf T}$ exchanges rows and columns. Consequently, $(AA^{\mathsf T})_{ij}$ is the dot product of row $i$ with row $j$. A matrix identity $AA^{\mathsf T}=nI$ says its rows have squared length $n$ and different rows are orthogonal. Here $I$ is the identity matrix, with ones on the diagonal and zeros elsewhere. This is the exact identity we will later read in a formalized research statement.

\begin{lab}{Replacing equation counting by a structural test}
\noindent\textbf{Deliberately flawed reply.} ``There are two linear equations and two unknowns, so the solution is unique.''

\noindent\textbf{Learner.} ``Test the matrix with rows $(1,1)$ and $(2,2)$. Find a nonzero vector it sends to zero. Then state a correct criterion for uniqueness.''

\noindent\textbf{Model answer.} The vector $(1,-1)$ is in the kernel. Any existing solution can be changed by a multiple of this vector without changing the outputs. Uniqueness of solutions, when they exist, follows from a zero kernel; existence for every target is a separate surjectivity claim. Ask the assistant to give one target that has solutions and one that does not for this matrix.
\end{lab}

\section*{Exercises}
\begin{exercise}\label{ex:10.1}For $A=\begin{pmatrix}2&1\\0&3\end{pmatrix}$ and $x=(1,-2)$, compute $Ax$. Then explain the answer using the columns of $A$.\end{exercise}
\begin{exercise}\label{ex:10.2}Find the kernel of the linear map $T:\RR^2\to\RR$ given by $T(x,y)=x+y$. Then describe all solutions of $T(x,y)=5$.\end{exercise}
\begin{exercise}\label{ex:10.3}Show that the map $T:\RR^2\to\RR^2$ given by $T(x,y)=(x+y,x-y)$ has inverse $(s,t)\mapsto\bigl((s+t)/2,(s-t)/2\bigr)$. Check both compositions.\end{exercise}
\begin{exercise}\label{ex:10.4}Compute the dot product and the lengths of $u=(1,2)$ and $v=(2,-1)$. Check the Cauchy--Schwarz inequality for this pair.\end{exercise}
\begin{exercise}\label{ex:10.5}Using the Cauchy--Schwarz inequality, prove the triangle inequality $\norm{u+v}_2\le\norm{u}_2+\norm{v}_2$ for all $u,v\in\RR^n$. Begin by comparing the squares of the two sides.\end{exercise}
\begin{exercise}\label{ex:10.6}Let $u_1,\ldots,u_k\in\RR^n$ be nonzero vectors that are pairwise orthogonal, which means that $u_i\cdot u_j=0$ whenever $i\ne j$. Show that $a_1u_1+\cdots+a_ku_k=0$ forces every coefficient $a_j$ to be zero; in other words, the list $u_1,\ldots,u_k$ is linearly independent. Then give an example showing that the word ``nonzero'' cannot be left out.\end{exercise}

